% TOP_PROBLEM: {"id": 3244, "title": "Acyclic list colouring of planar graphs.", "category_id": 3, "category": {"id": 3, "name": "graph_theory", "display_name": "Graph Theory", "description": "Problems involving graphs, networks, and their properties.", "slug": "graph-theory", "order_index": 3, "created_at": "2026-07-31T15:26:25.671Z"}, "status": "open", "problem_number": "OPG-46940", "set_id": 12}
% FIRST_POSTED: 2026-10-01
% ENTRY_KIND: statement_only
% UnsolvedMath ID 3244; source catalogue code OPG-46940
% !TeX program = lualatex
% Definition and sources only; this file is not an LLM solution attempt.
\documentclass[11pt,a4paper]{article}
\usepackage[margin=25mm]{geometry}
\usepackage{fontspec}
\setmainfont{lmroman10-regular.otf}[BoldFont=lmroman10-bold.otf,
 ItalicFont=lmroman10-italic.otf,BoldItalicFont=lmroman10-bolditalic.otf]
\usepackage{amsmath,amssymb,mathtools}
\usepackage{unicode-math}
\setmathfont{latinmodern-math.otf}
\AtBeginDocument{\let\setminus\smallsetminus}
\usepackage{xurl,hyperref}
\hypersetup{colorlinks=true,urlcolor=blue}
\setcounter{secnumdepth}{0}
\setlength{\parindent}{0pt}
\setlength{\parskip}{5pt}
\setlength{\emergencystretch}{3em}
\begin{document}
\section*{Acyclic list colouring of planar graphs.}\label{3244}
{\small UnsolvedMath ID 3244 · OPG-46940 · Graph Theory · OpenGarden}
\subsection{Definitions and mathematical statement}
Let $G=(V,E)$ be a finite simple planar graph. A list assignment gives each
vertex $v$ a finite set $L(v)$ of permitted colours. An $L$-colouring is a
map $c:V\to\bigcup_vL(v)$ satisfying $c(v)\in L(v)$ for every vertex.
It is proper when adjacent vertices receive different colours.

A proper colouring is acyclic if there is no cycle using only two colours.
Equivalently, for every two distinct colour names $a,b$, the induced
subgraph on $c^{-1}(\{a,b\})$ is a forest, meaning that it has no
undirected cycle. This condition permits cycles using three or more
colours. Say that $G$ is acyclically $k$-choosable if every assignment with
$|L(v)|=k$ for all $v$ admits an acyclic proper $L$-colouring.

The conjecture is that every finite simple planar graph is acyclically
five-choosable. Written with all quantifiers, for every such $G$ and every
family of five-element lists $\{L(v):v\in V\}$, there exists a map $c$
from those lists which is proper and for which each cycle has at least
three distinct colours among its vertices.

Lists are arbitrary and may overlap in any way; their union need not
have size five. Thus a colouring from one common five-colour palette
would establish only the ordinary acyclic-colouring assertion. Likewise,
a proper list-colouring without the condition on two-coloured cycles
would not establish this conjecture. Lists larger than five are equivalent
for this purpose, since one may first retain any five colours per vertex.
\subsection{Short English statement}
Can every planar graph be coloured from arbitrary lists of five colours so that every cycle uses at least three colours?
\subsection{Sources}
\begin{itemize}
\item[S1] ULAM.AI and the UnsolvedMath Contributors, supplied problems.json, record 3244 (OPG-46940), OpenGarden collection. Catalogue identity and source transcription; CC BY 4.0.
\url{https://huggingface.co/datasets/ulamai/UnsolvedMath}.
\item[S2] Open Problem Garden, Acyclic list colouring of planar graphs.. Original problem and discussion; the mathematical formulation below is expanded from the supplied source record.
\url{http://www.openproblemgarden.org/op/acyclic_list_colouring_of_planar_graphs}.
\end{itemize}
\end{document}
